\section*{第 \ref{ch:b2:living-soil} 章——有生命的土壤}

\begin{solution}{exo:b2:living-soil:1}
O 层：凋落物和部分分解的有机物，分解
始于此。A 层：与\omterm{def:b2:living-soil:soil}{腐殖质}混合的矿质\omterm{def:b2:living-soil:soil}{土壤}，颜色深暗，呈团粒结构，
容纳大部分根系、生物体、养分和水分。B 层：
淀积层，淋溶自上方的黏粒、铁氧化物和碳酸盐
在此沉积。C 层：风化母质，岩石
碎屑处于细粒基质中，新的矿质\omterm{def:b2:living-soil:soil}{土壤}在此释放。R 层：
未风化的基岩。
\end{solution}

\begin{solution}{exo:b2:living-soil:2}
黏粒与\omterm{def:b2:living-soil:soil}{腐殖质}表面的负电荷吸附可交换的
阳离子，根系可释放质子与之交换而吸收。砂粒
为石英，不带电且颗粒粗，单位质量的表面积小，
几乎不含\omterm{def:b2:living-soil:soil}{腐殖质}。施石灰即加入碳酸钙：钙
置换出使酸性\omterm{def:b2:living-soil:soil}{土壤}交换位点饱和的氢离子和铝离子，
碳酸根中和酸度，
交换位点随即被一种营养阳离子重新占据。
\end{solution}

\begin{solution}{exo:b2:living-soil:3}
分解者：细菌（\emph{Bacillus}）和真菌
（\emph{Trichoderma}）；微生物取食者：\omterm{def:b2:unicellular-diversity:protist}{原生生物}（变形虫）和
以细菌为食的线虫；\omterm{prop:b2:living-soil:community}{食碎屑者}：蚯蚓、马陆、
跳虫、潮虫；捕食者：螨类、捕食性线虫、
蜈蚣、步甲；顶层为鼹鼠或尖鼠。
取食者的 $\mathrm{C:N}$ 高于其所食细菌，
以铵的形式排出多余的氮：它们矿化了
细菌所固定的氮。
\end{solution}

\begin{solution}{exo:b2:living-soil:4}
丛枝菌根：真菌（球囊菌）进入根皮层
细胞，形成丛枝；见于大多数草本植物和作物以及
热带树木。外生菌根：真菌（担子菌和
子囊菌，其中许多为蘑菇）包裹根尖，并
在细胞之间生长而不进入细胞；见于温带
及寒带树木。两者中，植物提供糖，真菌
提供磷酸盐、氮、水、微量元素以及一些
对病原体的防护。
\end{solution}

\begin{solution}{exo:b2:living-soil:5}
$L^{*} = I/k = 4/0.8 = \qty{5}{t/ha}$；\omterm{def:b2:biogeochemical-cycles:reservoir}{停留时间} $1/k = 1.25$
年；\qty{90}{\%} 的损失需 $\ln 10/0.8 = 2.9$ 年。
\end{solution}

\begin{solution}{exo:b2:living-soil:6}
每摄取 \qty{100}{kg} 碳，微生物保留 \qty{40}{kg}，并
需要 $40/8 = \qty{5}{kg}$ 的氮。秸秆提供 $100/80 =
\qty{1.25}{kg}$：每 \qty{100}{kg}
碳固定 \qty{3.75}{kg}。三叶草提供 $100/15 = \qty{6.7}{kg}$：矿化
\qty{1.7}{kg}。
\end{solution}

\begin{solution}{exo:b2:living-soil:7}
经蚯蚓每年搬运 \qty{20}{t/ha} \omterm{def:b2:living-soil:soil}{土壤}。顶部
\qty{20}{cm} 厚的\omterm{def:b2:living-soil:soil}{土层}重 $0.2\times 10^{4}\times 1.3 = \qty{2600}{t/ha}$：
需 130 年才经过一次。达尔文的每英亩十吨
即 \qty{25}{t/ha}，量级相同——他所说的“每隔几年”指的是
表层的虫粪，厚仅几厘米，而非整个表土。
\end{solution}

\begin{solution}{exo:b2:living-soil:8}
细菌：$10^{9}\times 10^{-12} = \qty{1}{mg}$ 每克。菌丝：
体积 $\pi(2\times 10^{-4})^{2}\times 10^{4} = \qty{1.26e-3}{cm^3}$，
质量 \qty{1.4}{mg} 每克。每公顷（$2.6\times 10^{9}$ g）：
细菌 \qty{2.6}{t}，真菌 \qty{3.6}{t}。
\end{solution}

\begin{solution}{exo:b2:living-soil:9}
$H^{*} = 1.5/0.02 = \qty{75}{t/ha}$ 碳。翻耕时：$1.5/0.04
= \qty{37.5}{t/ha}$；损失 \qty{37.5}{t/ha}，其中一半
在 $\ln 2/0.04 = 17$ 年后流失。
\end{solution}

\begin{solution}{exo:b2:living-soil:10}
磷酸盐：$\sqrt{2\times 10^{-13}\times 8.64\times 10^{6}} =
\qty{1.3}{mm}$ 一季内——相距 \qty{1}{cm} 的根
只能触及两根之间\omterm{def:b2:living-soil:soil}{土壤}的八分之一，而相距 \qty{0.1}{mm} 的菌丝
可触及其中全部\omterm{def:b2:living-soil:soil}{土壤}。硝酸盐：\qty{4}{cm}，远超
根间距：它依靠自身（以及随水分的质流）到达根部，
因此菌根对硝酸盐作用甚微，
而对磷酸盐作用显著。
\end{solution}

\begin{solution}{exo:b2:living-soil:11}
质量 $0.25\times 10^{4}\times 1.3 = \qty{3250}{t/ha}$。净损失
\qty{19}{t/ha}（每年）：寿命 170 年。碳损失 $20\times
0.03 = \qty{0.6}{t/ha}$ 每年。免耕时净损失为
\qty{1}{t/ha}：寿命 3250 年。
\end{solution}

\begin{solution}{exo:b2:living-soil:12}
凋落物一两年周转一次，\omterm{def:b2:living-soil:soil}{腐殖质}约五十年，矿质 A
层以千年计（其形成速率为每年 \qty{0.1}{mm}），而
盐分在数十年内积累，若存在
排水条件，则亦在数十年内排出。农民一季内即可看到凋落物与作物的反应，
而\omterm{def:b2:living-soil:soil}{腐殖质}则在一生之内可见（下降三分之一需
三十年，恢复同样需要三十年）；矿质\omterm{def:b2:living-soil:soil}{土壤}的侵蚀
一生之内不可见，十年内即不可逆；盐分
一代之内可见，若无排水则永久存在。\omterm{def:b2:living-soil:soil}{土壤}的
缓变量，正是年度账目所不记录的。
\end{solution}

\begin{solution}{pb:b2:living-soil:1}
\textbf{1.} $0.25\times 10^{4}\times 1.3 = \qty{3250}{t/ha}$;
碳 $\qty{81}{t/ha}$。
\textbf{2.} $81/12 = \qty{6.8}{t/ha}$ 的氮。
\textbf{3.} 细菌 \qty{1}{mg/g}，$\qty{3.25}{t/ha}$；菌丝
$\pi(2\times 10^{-4})^{2}\times 2\times 10^{4}\times 1.1 =
\qty{2.8}{mg/g}$，\qty{9}{t/ha}：微生物生物量约
\qty{12}{t/ha}.
\textbf{4.} 活体物质约 \qty{14}{t/ha}（鲜重），
约占\omterm{def:b2:living-soil:soil}{腐殖质}碳的 \qty{17}{\%}——以碳计仅占其中几个百分点，
因为生物体主要是水。
\textbf{5.} 微生物碳约为 \qty{12}{t} 的一半，即 \qty{6}{t}；
每年周转两次，保留 \qty{40}{\%}，则其消耗
$2\times 6/0.4 = \qty{30}{t}$ 碳，并呼吸释放 \qty{18}{t/ha}
碳，相当于 \qty{66}{t} $\mathrm{CO_2}$（粗略估计：
保留率与周转率均为粗略值）。
\textbf{6.} 估计值为作物净初级生产量的三倍，
\qty{5}{t}：估计偏高，因为所作假设（全部生物量每年周转两次、
且无\omterm{def:b2:angiosperm-reproduction:seed}{休眠}部分）过于宽松；
作物下的实际\omterm{prop:b2:living-soil:humus}{土壤呼吸}与作物自身产量同一量级，
即每公顷每年数吨碳。\omterm{def:b2:living-soil:soil}{土壤}的
呼吸量与其上方的植物相当。
\textbf{7.} 碳 $5\times 0.45 = \qty{2.25}{t}$；氮 $5\times
0.03 = \qty{150}{kg}$; $\mathrm{C:N} = 15$.
\textbf{8.} 矿化：微生物需要 $2250\times 0.4/8 =
\qty{112}{kg}$ 氮，而残茬中含 150，因此
净释放 \qty{38}{kg/ha}。
\textbf{9.} $1 - e^{-0.5} = \qty{39}{\%}$（3 个月后）；$1 - e^{-2}
= \qty{86}{\%}$（一年后）。
\textbf{10.} $0.39\times 38 = \qty{15}{kg/ha}$.
\textbf{11.} 作物需要 \qty{150}{kg}，其中大部分在最初
两个月内需要；残茬在三个月内提供 \qty{15}{kg}，并
共计 \qty{38}{kg}——是补充，而非供给。其价值
还在于三叶草所固定、并保存于残茬中的氮，
即残茬自身的 \qty{150}{kg}，其中大部分进入\omterm{def:b2:living-soil:soil}{腐殖质}，
并在数年间陆续回归。
\textbf{12.} 秸秆：碳 \qty{2250}{kg}，氮 $2250/80 =
\qty{28}{kg}$；微生物需要 \qty{112}{kg}，因此 \qty{84}{kg/ha}
从\omterm{def:b2:living-soil:soil}{土壤}中固定——超过作物需求量的一半。
农民必须随秸秆施入氮肥，或在播种前将其充分翻耕入土，
或将其留在地表，使其缓慢腐解。
\textbf{13.} 翻耕：$1.2/0.025 = \qty{48}{t/ha}$；免耕：
$1.2/0.015 = \qty{80}{t/ha}$.
\textbf{14.} $H(t) = 80 - 32\,e^{-0.015t}$：增益 \qty{4.5}{t/ha}
（10 年后），\qty{17}{t/ha}（50 年后）。
\textbf{15.} $\mathrm{d}H/\mathrm{d}t = 0.015\times 32 =
\qty{0.48}{t/ha}$ 碳（第一年），相当于 \qty{1.8}{t} 的
$\mathrm{CO_2}$.
\textbf{16.} 第一年内 $0.48\times 1.5\times 10^{9} = \qty{0.7}{GtC}$，
占化石燃料排放的 \qty{7}{\%}；到第五十年，
速率降至 $0.48\,e^{-0.75} = \qty{0.23}{t/ha}$，
\qty{0.34}{GtC}.
\textbf{17.} 增益是趋近新\omterm{def:b2:biogeochemical-cycles:reservoir}{稳态}的过程：它
随 $H$ 达到 \qty{80}{t/ha} 而衰减为零，整个增益
若农田再次翻耕，则会在数十年内返回大气
——\omterm{def:b2:living-soil:soil}{土壤}\omterm{prop:b2:biogeochemical-cycles:carbon}{碳汇}是一次性的可逆转移，而非持久
对持续排放的抵消。
\textbf{18.} $k_h = 0.02$：$H^{*} = \qty{60}{t/ha}$，损失
\qty{20}{t/ha}，超过前五十年免耕的增益：
变暖可能抵消
这一管理措施的效果。
\textbf{19.} 翻耕：$15 - 0.5 = \qty{14.5}{t/ha}$ 每年；免耕：
\qty{1}{t/ha}.
\textbf{20.} $3250/14.5 = 220$ 年；$3250/1 = 3250$ 年。
\textbf{21.} 碳 $15\times 0.025 = \qty{0.38}{t/ha}$；氮
$0.38/12 = \qty{31}{kg/ha}$ 每年。
\textbf{22.} 翻耕\omterm{def:b2:biogeochemical-cycles:reservoir}{稳态}下\omterm{def:b2:living-soil:soil}{腐殖质}分解量为 $k_hH^{*} =
I = \qty{1.2}{t/ha}$ 每年；侵蚀又增加其三分之一，
但与分解不同，它不能由输入平衡——它
使储量耗尽。
\textbf{23.} 部分如此。被侵蚀的碳埋藏于沉积物或
\omterm{def:b2:biogeochemical-cycles:reservoir}{储库}中，可长期免受分解（埋藏
型\omterm{prop:b2:biogeochemical-cycles:carbon}{碳汇}）；但团聚体在搬运中被破坏，大部分
碳在途中被氧化，而农田流失的肥力
由化肥补充，而化肥生产会排放碳。全球
侵蚀所致碳效应的全球净方向仍有争议。
\textbf{24.} 凋落物 $1/k = 0.5$ 年；\omterm{def:b2:living-soil:soil}{腐殖质} $1/k_h = 40$ 年
（翻耕时）；免耕时为 67 年；矿质\omterm{def:b2:living-soil:soil}{土壤} $3250/0.5 =
6500$ 年（就形成而言），而就
流失而言，翻耕下为 220 年。
\textbf{25.} A 层碳储量 \qty{81}{t/ha}；寿命
翻耕下 220 年，免耕下 3250 年；三叶草残茬
约提供 \qty{38}{kg/ha} 矿质氮，其中 15 在
最初三个月内释放。
\end{solution}
